\section{Decision-scale numerical and statistical precision}\label{sec:precision99}
The continuous qualification uses the original terminal-controller allowance and all-action lower witness, not just a small surrogate residual. Convex midpoint and endpoint risk rules supply a width of at most $KR/\{(1-\alpha)n^2\}$ with $2n+1$ continuation queries. Theorem~\ref{thm:nonuniform99} below extends this account beyond the uniform law. The full uniform-rule statement, numerical endpoint charges and tensor extension remain in Section~\ref{sec:precision96} of the Development Appendix.

For the frozen population comparison, every complete service is qualified to the original-optimum allowance $\tau=1/256$. The design uses 4,096 independent initial states at each risk level and twelve simultaneous contrasts, with numerical half-width and sampling allowance totaling $3/8192$. Theorem~\ref{thm:power96} states the exact simultaneous event and its independence requirements. Those population intervals belong to the frozen actors, not to subsequently refitted candidates. The acquisition comparison instead retains the deterministic contrast support $[-\tau,\tau]$.

For a workload of $B$ decisions, the purchase account combines the per-decision policy contrast, $B$ times that contrast, full-workload computation and one-time acquisition. Theorems~\ref{thm:purchase97} and Proposition~\ref{prop:purchase96} retain the complete interval and strict-boundary statements. Historical timing ranges are descriptive inputs to a conditional price exercise. The exact decision mechanism in Section~\ref{sec:decision99} will provide an identified model-specific policy contrast; it will not turn these historical clock ranges into confidence intervals or external prices into estimated welfare weights.
